\subsection*{A faithful smaller ambient group from the common involution}

\statement{B.73}{Proved}{The common-involution quotient preserves the full completed intersection}
Use B.65's actual involution
\[
 r=(1\ 3)(2\ 5)(6\ 8)(12\ 10)(9\ 11)(4\ 7)
\]
and put \(C_r=C_{P\Gamma_{0,12}}(r)\).
Let \(Q_8\) be the mapping-class group of its quotient sphere,
fixing the six ordinary orbit labels individually and preserving
the unordered pair of branch labels.
There is an explicitly specified subgroup \(L_8\leq Q_8\) of index
\(32\) and a natural isomorphism
\[
 C_{\widehat{P\Gamma}_{0,12}}(r)
       =\widehat{C_r}\ \xrightarrow{\ \sim\ }\widehat{L_8}
                  \ \lhook\joinrel\longrightarrow\widehat{Q_8}.
                                                               \tag{B.73.1}
\]
For the actual source images \(A,B\) of \(H_T,H_5\), respectively,
this identifies
\[
             \mathcal K\ \xrightarrow{\ \sim\ }
                      \overline A\cap\overline B
                         \quad\hbox{in }\widehat{Q_8}.           \tag{B.73.2}
\]
In addition, forgetting the ordinary label from \(\{1,3\}\)
gives a map to the analogous colored seven-point group \(Q_7\).
Its restriction to the actual \(\widehat H_5\), and hence to
\(\mathcal K\), is injective.  The image of \(\mathcal K\)
lies in the open pure subgroup \(\widehat{P\Gamma}_{0,7}\).
No equality between \(\mathcal K\) and the intersection of
the two seven-point source images is asserted.

\begin{proof}
The branch-ordered comparisons rederive LM4 F.12 and F.14.
We give the full colored-source formulation, its sheet indices,
and the specified lift sections explicitly.

\emph{The actual eight-point group.}
All twelve original marks occur in free \(r\)-orbits.
The quotient has six ordinary marks and the images of the
two unmarked fixed points of \(r\).
Every class in \(Q_8\) lifts through the double cover:
its monodromy character has value one at each branch
meridian and zero at every ordinary meridian, and \(Q_8\)
preserves this character.
The regular fully ramified covering comparison used in
B.28--B.30 therefore identifies the symmetric mapping-class
group preserving each ordinary orbit, modulo its deck
involution, with \(Q_8\).

Changing a lift by \(r\) changes its sheet permutation
above all six ordinary marks simultaneously.
The sheet bits consequently give a well-defined
homomorphism
\[
 \epsilon_8:Q_8\longrightarrow
       (\mathbb Z/2)^6/\langle(1,1,1,1,1,1)\rangle.              \tag{B.73.3}
\]
A point push of one ordinary mark around one branch mark
realizes its single coordinate bit, so this map is onto.
Define \(L_8=\ker(\epsilon_8)\).
Exactly one of the two lifts of a class in \(L_8\) fixes
all twelve original labels.  Thus the quotient comparison
identifies \(C_r\) with \(L_8\), and
\([Q_8:L_8]=32\).

The choice of lift can also be made explicitly on all of
\(Q_8\): choose the ordinary mark from \(\{1,3\}\), and
require the lift to fix original label \(1\).
Exactly one lift does so, and this condition is preserved
under composition.  It defines a splitting of the central
deck extension.  On \(L_8\) it is the all-label-pure lift
just described.

The prime-order-two centralizer comparison and full-topology
argument in B.30 apply to this particular \(r\).
They identify \(C_{\widehat{P\Gamma}_{0,12}}(r)\) with
the full completion of \(C_r\), as its actual embedded
closure.  Since \(L_8\) has finite index in \(Q_8\), its
completion is the corresponding open subgroup of
\(\widehat{Q_8}\).  This proves (B.73.1).
B.65 puts both actual source groups in \(C_r\).
Their completions lie there by continuity, so the injective
comparison preserves their intersection and proves
(B.73.2).

\emph{The seven-point map and its faithful cyclic restriction.}
First forget original labels \(1,3\), the two \(D\)-fixed
poles.  Write \(D_{10}\) and \(r_{10}\) for the resulting
order-five and order-two mapping classes on the ten
retained marks.  The restricted forgetting homomorphism is
an isomorphism
\[
 H_5=C_{P\Gamma_{0,12}}(D)
       \ \xrightarrow{\ \sim\ }\
       C_{P\Gamma_{0,10}}(D_{10}).                              \tag{B.73.4}
\]
To see this without an isotopy ambiguity, use the cyclic
covering comparison on both sides.  Both quotient orbifolds
have the same two ordinary orbit marks and two branch
points.  The latter are individually distinguished by
the local rotation residues \(+1,-1\pmod5\).
A representative commuting with \(D_{10}\) preserves its
two fixed points individually; marking them recovers
the inverse comparison.
The deck rotations act nontrivially on the ten retained
labels, so there is no pure finite kernel on either side.
The resulting pure group is the same rank-six group of
the parameter curve
\(\mathbb P^1_u\setminus\{0,\infty,\mu_5\}\).
This is also the direct quotient-cover identification
in B.28.

The valid prime-order-five full-topology and density theorem
applies to both surfaces in (B.73.4).
Their completed centralizers are therefore the completions
of those same discrete groups, compatibly with forgetting.
It follows that the restricted map on the actual
\(\widehat H_5\) is injective.
Discrete injectivity alone is not being used to infer
this completed assertion.

The involution \(r_{10}\) is free on the ten labels.
Its quotient has five ordinary labels and an unordered
branch pair.  The same sheet-bit construction gives
\[
 \epsilon_7:Q_7\longrightarrow
       (\mathbb Z/2)^5/\langle(1,1,1,1,1)\rangle,\qquad
 L_7=\ker(\epsilon_7),\qquad [Q_7:L_7]=16,                       \tag{B.73.5}
\]
and
\[
 C_{\widehat{P\Gamma}_{0,10}}(r_{10})
          \simeq\widehat{L_7}
                 \lhook\joinrel\longrightarrow\widehat{Q_7}.
                                                               \tag{B.73.6}
\]
Here the chosen section fixes original label \(2\) above
the ordinary mark from \(\{2,5\}\).
The valid prime-order-two comparison proves (B.73.6).
Compose (B.73.4) after completion with (B.73.6).
This proves the claimed faithful seven-point comparison
on \(\widehat H_5\).
The two chosen sections need not commute under forgetting
on every class of \(Q_8\); their quotient maps and their
all-label-pure restrictions do commute.

\emph{Coordinates and the pure target on \(\mathcal K\).}
In the cyclic family \(r_u(z)=u/z\), the quotient coordinate
and ordinary marks are
\[
 t=z+u/z,\qquad U=\infty,\qquad
 t_i(u)=\zeta^i+u\zeta^{-i}\quad(0\leq i<5),\qquad
 \zeta^5=1,\quad\zeta\ne1.                                    \tag{B.73.7}
\]
The branch pair is \(\{2\sqrt u,-2\sqrt u\}\).
The seven-point map forgets \(U\) and keeps both branch
marks.  A loop about \(u=0\) exchanges them, so the whole
cyclic source must land in the colored group \(Q_7\).
On the index-two source \(u=s^2\), normalize
\(-2s,2s,t_0\) to \(0,\infty,1\).
The four remaining coordinates are
\[
       \left(
       \frac{(s+\zeta^i)(s-1)}
            {(s-\zeta^i)(s+1)}
       \right)^2,\qquad 1\leq i\leq4.                          \tag{B.73.8}
\]
This ordered source is
\(\mathbb P^1_s\setminus\{0,\infty,\mu_{10}\}\), of free
rank eleven.
The based marking is obtained from the actual B.61/B.65
common marking by the displayed quotient and forgetting;
no independent ambient conjugator is introduced.

For the ordered tetrahedral family the actual \(T\)-action
is fixed throughout.  Its equivariant ambient isotopies
can retain the two fixed points of the chosen \(r\)
individually.  Its quotient monodromy consequently has
trivial branch-pair permutation, and the same holds on
\(\widehat H_T\) by continuity.
Thus the image of \(\mathcal K\) in \(Q_7\) has trivial
branch permutation and lies in the stated pure completion.
These two extra branch marks were not original punctures;
they were absent from B.69's rigid tests.
The faithful map here therefore does not contradict
that actual full-rigid relation.
\end{proof}

The eight-point comparison also transfers B.72:
\[
                    A\cap qBq^{-1}=1
                              \qquad(q\in Q_8).                \tag{B.73.9}
\]
Indeed lift \(q\) to an actual twelve-point mapping class.
An equality between the two pure source lifts modulo the
deck group has trivial deck difference, because \(r\)
is not pure.  B.72 then applies.
The unrefined product \(AB\) is closed in \(Q_8\) by B.52
and the full topology of \(C_r\simeq L_8\).
Thus B.46 gives the exact remaining quantifier
\[
 \mathcal K=1\quad\Longleftrightarrow\quad
 (A\cap V)B\text{ is closed in }Q_8
                    \quad\text{for every }V\leq_f Q_8.         \tag{B.73.10}
\]
This condition remains unresolved.
In particular neither (B.73.9) nor the individual source
embeddings provide a simultaneous virtual retraction.

